\documentclass[11pt]{ctexart}

\usepackage[a4paper,margin=1in]{geometry}
\usepackage{amsmath,amssymb}
\usepackage{enumitem}
\usepackage{xcolor}
\usepackage{hyperref}
\usepackage{fancyhdr}

\hypersetup{
  colorlinks=true,
  linkcolor=black,
  urlcolor=blue
}

\definecolor{agenblue}{RGB}{28,83,168}
\definecolor{agengray}{RGB}{90,96,110}

\setlength{\parindent}{0pt}
\setlength{\parskip}{0.58em}
\setlength{\headheight}{14pt}
\setlist[itemize]{leftmargin=1.5em,itemsep=0.25em,topsep=0.25em}
\setlist[enumerate]{leftmargin=1.6em,itemsep=0.25em,topsep=0.25em}
\allowdisplaybreaks
\sloppy

\pagestyle{fancy}
\fancyhf{}
\lhead{\textcolor{agengray}{AgenQA Pre-Release}}
\rhead{\textcolor{agengray}{Representative Path View Samples}}
\cfoot{\thepage}
\renewcommand{\headrulewidth}{0.4pt}

\newcommand{\sampletitle}[2]{
  \section*{\textcolor{agenblue}{#1}}
  \textbf{Subject:} #2
  \vspace{0.25em}
  \hrule
  \vspace{0.6em}
}

\newcommand{\focus}[1]{\textbf{Path View focus.} #1}
\newcommand{\why}[1]{\textbf{Why this sample matters.} #1}

\begin{document}

{\LARGE \textbf{AgenQA 代表性 Path View 样例题}}

\vspace{0.4em}

{\large Selected solver-facing questions for the AgenQA pre-release showcase}

\vspace{0.8em}

本文展示三道 representative \textbf{Path View} sample questions，用于说明 AgenQA 生成题目的形态与难度来源。它们不是完整 benchmark release，也不附带 source papers、raw solver responses、prompt snapshots 或 run artifacts。

这些样例都属于 solver-facing Path View：题面只给出最终求解所需的可见条件，中间 dependency path 被折叠，解题者需要自行恢复关键推理步骤。

\sampletitle{Sample 1. FedAvg 收敛界中的学习率选择}{Federated and Distributed Learning}

\focus{从 FedAvg 的 smoothness、non-IID 梯度异质性和通信轮数条件出发，恢复收敛上界中的学习率优化步骤。}

\textbf{Solver-facing question.}

给定 FedAvg 设定：
\begin{itemize}
  \item 全局目标函数为
  \[
  F(w)=\sum_k \frac{n_k}{n}F_k(w),
  \]
  其中 $F_k(w)$ 是客户端 $k$ 的局部经验损失，$n_k=|D_k|$，且 $n=\sum_k n_k$。
  \item 每个客户端从 $w^t$ 出发，以常数学习率 $\eta>0$ 执行 $E>1$ 步本地 SGD，再由服务器按权重平均聚合。
  \item 客户端数据 non-IID，梯度异质性由加权 RMS 量 $G$ 表示：
  \[
  G=\left(\sum_k \frac{n_k}{n}\|\nabla F_k(w)-\nabla F(w)\|^2\right)^{1/2}\geq 0.
  \]
  \item 全局损失满足 $L$-smoothness：
  \[
  F(w')\leq F(w)+\langle \nabla F(w),w'-w\rangle+\frac{L}{2}\|w'-w\|^2.
  \]
  \item 一共进行 $T$ 轮全客户端参与通信，$E,G,L,T$ 均为正数。
\end{itemize}

要求求出使时间平均平方梯度范数上界最小的学习率 $\eta^\ast$，并给出优化后的最紧上界。最终答案只能使用 $F(w^0)-F(w^\ast)$、$L$、$E$、$G$、$T$。

\textbf{Reference answer.}
\[
\boxed{
\eta^\ast=
\sqrt{\frac{2[F(w^0)-F(w^\ast)]}{L E^2G^2T}}
}
\]
\[
\boxed{
\frac{1}{T}\sum_{t=0}^{T-1}\|\nabla F(w^t)\|^2
\leq
2\sqrt{\frac{2LG^2[F(w^0)-F(w^\ast)]}{T}}
}
\]

\why{题面不直接给出中间收敛界的单步优化式，而要求 solver 从 FedAvg 设定中恢复出学习率选择与优化后上界。}

\sampletitle{Sample 2. Pointer-Generator Coverage Loss 的梯度路径}{Natural Language Processing / Abstractive Summarization}

\focus{在 pointer-generator + coverage loss 结构中，恢复编码器隐藏状态 $h_i$ 到总损失的两条梯度路径。}

\textbf{Solver-facing question.}

给定软注意力 encoder-decoder：
\[
e_{t,i}=\mathrm{score}(s_t,h_i),\qquad
a_{t,i}=\frac{\exp(e_{t,i})}{\sum_j\exp(e_{t,j})},\qquad
c_t=\sum_i a_{t,i}h_i.
\]

组合训练损失为
\[
L_{\mathrm{total}}=L_{\mathrm{NLL}}+\lambda L_{\mathrm{cov}},
\]
其中
\[
\mathrm{cov}_{t,i}=\sum_{t'<t}a_{t',i},\qquad
L_{\mathrm{cov}}=\sum_t\sum_i\min(a_{t,i},\mathrm{cov}_{t,i}).
\]

编码器隐藏状态 $h_i$ 是一个向量，$\lambda\geq0$ 是固定标量，$\partial e_{t,i}/\partial h_i$ 表示标量 score 关于 $h_i$ 的向量值 Jacobian。由于 score 函数未固定，所以保持为未展开形式。

要求推导 $\partial L_{\mathrm{total}}/\partial h_i$ 的单一闭式符号表达。答案应同时捕捉：
\begin{itemize}
  \item $h_i\rightarrow e_{t,i}\rightarrow a_{t,*}$ 的 score path；
  \item $h_i\rightarrow c_t\rightarrow L_{\mathrm{NLL}}$ 的 context path。
\end{itemize}

\textbf{Reference answer.}
\[
\boxed{
\frac{\partial L_{\text{total}}}{\partial h_i}
=
\sum_t
\left[
a_{t,i}\left(g_{t,i}-\sum_j a_{t,j}g_{t,j}\right)
\frac{\partial e_{t,i}}{\partial h_i}
+
a_{t,i}\frac{\partial L_{\text{NLL}}}{\partial c_t}
\right]
}
\]

\why{题目要求 solver 把 softmax Jacobian、coverage dependency 和 context-vector path 合并，而不是只回答一个局部梯度公式。}

\sampletitle{Sample 3. TRADES 鲁棒训练目标的外层与内层条件}{Machine Learning Security / Adversarial Robustness}

\focus{从 TRADES-style robust objective 出发，同时恢复外层 envelope-theorem gradient 与内层 PGD / KKT 条件。}

\textbf{Solver-facing question.}

给定参数化深度神经网络 $f_\theta$，可行扰动集合为闭 $L_p$ 球：
\[
\{\delta:\|\delta\|_p\leq\varepsilon\}.
\]

TRADES-style 目标为：
\[
\min_{\theta}
\mathbb{E}_{(x,y)\sim\mu}
\left[
L(f_\theta(x),y)
+
\beta\max_{\|\delta\|_p\leq\varepsilon}
\mathrm{KL}\bigl(f_\theta(x)\|f_\theta(x+\delta)\bigr)
\right].
\]

其中 $\mathrm{KL}(f_\theta(x)\|f_\theta(x+\delta))$ 以 clean input distribution 为第一参数，以 perturbed input distribution 为第二参数。$\delta^\ast(\theta)$ 表示内层 KL 最大化问题的解。

要求合并推导：
\begin{enumerate}
  \item $\nabla_\theta\max_\delta\mathrm{KL}(\cdot)$ 的 envelope-theorem 形式，以及 $\theta^\ast$ 处外层驻点条件；
  \item 内层关于 $\delta$ 的归一化 PGD 更新规则；
  \item 收敛时的 KKT 条件，包括梯度驻点、互补松弛和可行性。
\end{enumerate}

\textbf{Reference answer.}
\[
\nabla_\theta
\max_{\|\delta\|_p\leq\varepsilon}
\mathrm{KL}\bigl(f_\theta(x)\|f_\theta(x+\delta)\bigr)
=
\nabla_\theta
\mathrm{KL}\bigl(f_\theta(x)\|f_\theta(x+\delta^\ast(\theta))\bigr).
\]
\[
\nabla_\theta\mathbb{E}_{(x,y)\sim\mu}
\left[L(f_{\theta^\ast}(x),y)\right]
+
\beta\nabla_\theta\mathbb{E}_{(x,y)\sim\mu}
\left[
\mathrm{KL}\bigl(f_{\theta^\ast}(x)\|f_{\theta^\ast}(x+\delta^\ast(\theta^\ast))\bigr)
\right]
=0.
\]
\[
\delta_{t+1}
=
\Pi_{\|\cdot\|_p\leq\varepsilon}
\left(
\delta_t+
\alpha
\frac{
\nabla_\delta\mathrm{KL}\bigl(f_\theta(x)\|f_\theta(x+\delta_t)\bigr)
}{
\left\|
\nabla_\delta\mathrm{KL}\bigl(f_\theta(x)\|f_\theta(x+\delta_t)\bigr)
\right\|_p
}
\right).
\]
\[
\nabla_\delta\mathrm{KL}\bigl(f_\theta(x)\|f_\theta(x+\delta^\ast)\bigr)
=
\lambda\nabla_\delta\|\delta^\ast\|_p,\qquad
\lambda\geq0.
\]
\[
\lambda(\|\delta^\ast\|_p-\varepsilon)=0,\qquad
\|\delta^\ast\|_p\leq\varepsilon.
\]

\why{题目覆盖 optimization / robustness reasoning：solver 需要同时处理外层参数优化、内层 adversarial maximization 和 constrained optimality conditions。}

\end{document}
